\documentclass[10pt,a4paper]{amsart}
\usepackage[margin=19mm]{geometry}
\usepackage{amsmath,amssymb,mathtools}
\usepackage[scr=rsfs,cal=euler]{mathalfa}
\usepackage{xfrac}
\usepackage[unicode,hidelinks,bookmarks=false]{hyperref}
\newcommand{\ie}{i.e.\ }
\newcommand{\cf}{cf.\ }
\newcommand{\resp}{respectively\ }
\newcommand{\Subsection}{\subsection}
\providecommand{\nobreakdash}{\nobreak\hbox{-}}
\setlength{\emergencystretch}{2em}
\multlinegap0pt
\usepackage{xcolor}
\usepackage{csquotes}
\usepackage{tikz}
\usepackage{amssymb}
\usepackage{dsfont}
\usepackage{mathtools}
\usepackage[shortlabels]{enumitem}
\usepackage{float}
\newtheorem{prop}{Proposition}
\newcommand{\End}{\setft{End}}
\newcommand{\Herm}{\setft{Herm}}
\newcommand{\I}{\mathds{1}}
\newcommand{\Sym}{\setft{Sym}}
\newcommand\V{\mathcal{V}}
\newcommand{\Vc}{\V_{\complex}}
\newcommand{\Vr}{\V_{\real}}
\def\ba#1\ea{\begin{align}#1\end{align}}
\newcommand{\bfalpha}{\boldsymbol{\alpha}}
\newcommand{\bfbeta}{\boldsymbol{\beta}}
\newcommand{\bff}{\mathbf{f}}
\newcommand{\bfgamma}{\boldsymbol{\gamma}}
\newcommand{\complex}{\mathbb{C}}
\newcommand{\dg}{\dagger}
\newcommand{\integer}{\mathbb{Z}}
\newcommand{\ip}[2]{\langle #1 , #2\rangle}
\newcommand{\norm}[1]{\lVert\tinyspace #1 \tinyspace\rVert}
\newcommand{\real}{\mathbb{R}}
\newcommand{\setft}[1]{\mathrm{#1}}
\renewcommand{\t}{{\scriptscriptstyle\mathsf{T}}}
\newcommand{\tinyspace}{\mspace{1mu}}
\title{\bf A hierarchy of eigencomputations for polynomial optimization on the sphere}
\author{Benjamin Lovitz, Nathaniel Johnston}
\date{Source: arXiv:2310.17827v3 (CC BY 4.0)}
\begin{document}
\maketitle

\noindent\emph{Source and licence.} \bf A hierarchy of eigencomputations for polynomial optimization on the sphere. Version arXiv:2310.17827v3 (CC BY 4.0), distributed by arXiv under the Creative Commons Attribution 4.0 International licence, http://creativecommons.org/licenses/by/4.0/, as recorded in the arXiv licence metadata for this exact version and shown on the abstract page https://arxiv.org/abs/2310.17827v3. Full licence notice: \texttt{licenses/arxiv-2310.17827-CC-BY-4.0.txt}.\par\smallskip

\noindent\textit{Editorial scope.} Counted unit: the article's prop eq:M (source0.tex lines 2655--2668). The article's complete original proof of it is delivered with it (source0.tex lines 2670--2714, one \texttt{proof} environment of 45 lines). No mathematics is written, repaired or re-derived: the statement and the proof are the article's own lines, in the article's own notation, and no retained line is substituted or re-flowed.  Retained uncounted as the source-local context the proof needs: lines 1258--1260.  Nothing was omitted from any retained span, so the package carries no omission record.  No retained line contains a citation, so the wrapper carries no bibliography and no bibliography entry.  No retained line contains a figure environment, an image-inclusion command or drawing code, so no image is delivered and the package declares no asset dependency.  Editorial formatting only, in addition: an AMS \texttt{amsart} preamble that restates the article's own declarations and macros used by the retained lines, namely the environments \texttt{prop} on separate counters, together with the standard packages those lines need.  The retained environments print in this document as Proposition 1 in order of appearance, whereas the article prints the same items with the numbering of its own full document.  The complete native source of the article as distributed in the arXiv e-print for this exact version is bundled unchanged as \texttt{sources/149141/source0.tex}.
\begin{prop}\label{prop:sym}
For a real form $p \in S^{2d}(\Vr^*)$ of even degree, $M(p) \in \Sym(S^{d}(\Vr))$ is the unique Gram operator for which $M(p)^{\t_1}=M(p)$.
\end{prop}

\begin{prop}
Let $p \in S^{2d}(\Vr^*)$ be a form, let $P=M(p) \in \Sym(S^d(\Vr))$ be the maximally symmetric Gram operator of $p$, and for $k \geq d$ let $P_k = M(p)_k \in \Sym(S^k(\Vr))$ be defined by $P_k= \Pi_k (P \otimes \I_{\Vr}^{\otimes k-d})\Pi_k$. If
\ba
p(x)=\sum_{\substack{\bfgamma \in \integer_{\geq 0}^n\\ |\bfgamma|=2d}} C_{\bfgamma} \binom{2d}{\bfgamma} x^{\bfgamma},
\ea
then
\ba\label{eq:M}
P=\sum_{|\bfalpha|=|\bfbeta|=d} C_{\bfalpha+\bfbeta} \binom{d}{\bfalpha} \binom{d}{\bfbeta} e^{\bfalpha}   (e^{\t})^{\bfbeta},
\ea
and
\ba\label{eq:Mk}
P_k=\sum_{\substack{|\bfalpha|=|\bfbeta|=d \\ |\bfgamma|=k-d}} C_{\bfalpha+\bfbeta} \binom{d}{\bfalpha} \binom{d}{\bfbeta} \binom{k-d}{\bfgamma} e^{\bfalpha+ \bfgamma}  (e^{\t})^{\bfbeta+\bfgamma}.
\ea
\end{prop}

\begin{proof}
Let $M \in \End(S^d(\Vr))$ be given by the righthand side of~\eqref{eq:M}. We use Proposition~\ref{prop:sym} to verify that $M=P$. First note that
\ba
\ip{x^{\otimes d}}{M x^{\otimes d}}&=\sum_{|\bfalpha|=|\bfbeta|=d} C_{\bfalpha+\bfbeta} \binom{d}{\bfalpha} \binom{d}{\bfbeta} x_1^{\alpha_1+\beta_1} \cdots x_n^{\alpha_n+\beta_n}\\
&=\sum_{\substack{\bfgamma \in \integer_{\geq 0}^n\\ |\bfgamma|=2d}} \sum_{\substack{|\bfalpha|=|\bfbeta|=d\\ \bfalpha+\bfbeta=\bfgamma}} C_{\bfgamma} \binom{d}{\bfalpha} \binom{d}{\bfbeta}  x_1^{\gamma_1} \cdots x_n^{\gamma_n}\\
&=\sum_{\substack{\bfgamma \in \integer_{\geq 0}^n\\ |\bfgamma|=2d}} C_{\bfgamma} \binom{2d}{\bfgamma} x_1^{\gamma_1} \cdots x_n^{\gamma_n}\\
&=p(x),
\ea
where the first two equalities are obvious, and the third follows from the identity
\ba
\sum_{\substack{|\bfalpha|=|\bfbeta|=d\\ \bfalpha+\bfbeta=\bfgamma}} \binom{d}{\bfalpha} \binom{d}{\bfbeta}=\binom{2d}{\bfgamma}.
\ea
It remains to prove that $M^{\t_1}=M$. Let $\bff_i \in \integer_{\geq 0}^n$ be the vector with $1$ in the $i$-th position and zeroes elsewhere, and note that $e^{\bfalpha}=\frac{1}{d} \sum_{i : \alpha_i >0} \alpha_i e_i \otimes e^{\bfalpha-\bff_i}$ when $|\bfalpha|=d$. Thus,
\ba
M&=\sum_{|\bfalpha|=|\bfbeta|=d} C_{\bfalpha+\bfbeta} \binom{d}{\bfalpha} \binom{d}{\bfbeta} e^{\bfalpha}   (e^{\t})^{\bfbeta}\\
&=\frac{1}{d^2}  \sum_{|\bfalpha|=|\bfbeta|=d} C_{\bfalpha+\bfbeta} \binom{d}{\bfalpha} \binom{d}{\bfbeta} \sum_{\substack{i,j \in [n]\\ \alpha_i, \beta_j >0}}  \alpha_i \beta_j (e_i \otimes e^{\bfalpha - \bff_i})  (e_j \otimes e^{\bfbeta - \bff_j})^\t.
\ea
Hence,
\ba
M^{\t_1}&=\frac{1}{d^2} \sum_{|\bfalpha|=|\bfbeta|=d} C_{\bfalpha+\bfbeta} \binom{d}{\bfalpha} \binom{d}{\bfbeta} \sum_{\substack{i,j \in [n]\\ \alpha_i, \beta_j >0}} \alpha_i \beta_j (e_j \otimes e^{\bfalpha - \bff_i})  (e_i \otimes e^{\bfbeta - \bff_j})^\t\\
&=\frac{1}{d^2} \sum_{i,j=1}^n  \sum_{\substack{|\bfalpha|=|\bfbeta|=d\\ \alpha_i, \beta_j>0}}C_{\bfalpha+\bfbeta}\binom{d}{\bfalpha} \binom{d}{\bfbeta}\alpha_i \beta_j (e_j \otimes e^{\bfalpha - \bff_i})  (e_i \otimes x^{\bfbeta - \bff_j})^\t\\
%&=\frac{1}{d^2} \sum_{i,j=1}^n  \sum_{\substack{|\bfalpha|=|\bfbeta|=d\\ \alpha_i, \beta_j>0}}  C_{\bfalpha+\bfbeta} \binom{d}{\bfalpha-\bff_i+\bff_j} \binom{d}{\bfbeta-\bff_j+\bff_i} (\alpha_j- \delta_{i,j}+1) (\beta_i-\delta_{i,j}+1) (x_j \otimes x^{\bfalpha - \bff_i})  (x_i \otimes x^{\bfbeta - \bff_j})^*\\
&= \frac{1}{d^2} \sum_{i,j=1}^n  \sum_{\substack{|\bfalpha'|=|\bfbeta'|=d\\ \alpha'_j, \beta'_i >0}}  C_{\bfalpha'+\bfbeta'} \binom{d}{\bfalpha'} \binom{d}{\bfbeta'} \alpha'_j \beta'_i (e_j \otimes e^{\bfalpha'-\bff_j})  (e_i \otimes e^{\bfbeta'-\bff_i})^\t\\
&=M,
\ea
where the third line follows from setting $\bfalpha'= \bfalpha - \bff_i + \bff_j$ and $\bfbeta'=\bfbeta-\bff_j+\bff_i$, and the rest are straightforward. Hence $M=P$.

Let $M_k$ be the righthand side of~\eqref{eq:Mk}. Let $\Vc=\complex^n$ (the complexification of $\Vr$), and regard $P$ as an element of $\Herm(S^d(\Vc))$ and $M_k$ as an element of $\Herm(S^k(\Vc))$. Recall that, as a Hermitian form, $P_k$ is given by
\ba
P_k(z,z^\dg)=P(z,z^\dg) \cdot \norm{z}^{2(k-d)}.
\ea
To complete the proof, it suffices to verify that $M_k(z,z^\dg)=P_k(z,z^\dg)$. Let
\ba
z = z_1 e_1 +\dots + z_n e_n \in \Vc.
\ea
Then
\ba
P_k (z,z^\dg) &=\sum_{|\bfalpha|=|\bfbeta|=d} C_{\bfalpha+\bfbeta} \binom{d}{\bfalpha} \binom{d}{\bfbeta} z^{\bfalpha}  (z^{\dg})^{\bfbeta} (z_1 z_1^\dg + \dots + z_n z_n^\dg)^{k-d}\\
&= \sum_{|\bfalpha|=|\bfbeta|=d} C_{\bfalpha+\bfbeta} \binom{d}{\bfalpha} \binom{d}{\bfbeta} z^{\bfalpha}  (z^{\dg})^{\bfbeta} \sum_{|\bfgamma|=k-d} \binom{k-d}{\bfgamma} z^{\bfgamma} (z^{\dg})^{\bfgamma} \\
&=\sum_{\substack{|\bfalpha|=|\bfbeta|=d \\ |\bfgamma|=k-d}} C_{\bfalpha+\bfbeta} \binom{d}{\bfalpha} \binom{d}{\bfbeta} \binom{k-d}{\bfgamma} z^{\bfalpha+ \bfgamma}  (z^{\dg})^{\bfbeta+\bfgamma}\\
&=\ip{z^{\otimes d}}{M_k z^{\otimes d}}\\
&=M_k(z,z^\dg).
\ea
This completes the proof.
\end{proof}

\end{document}
